%===============================================================================
% Research CV (academic version) — layout modelled on Matteo Ficarra's CV.
% This is the version published on guidoardizzone.github.io (files/cv.pdf),
% so it is web-safe: no phone number, no referee emails. Scholarship amounts are
% shown by choice (Sep 2026).
% Private-sector version: main.tex. World Bank YPP version: main_wbg.tex.
% To publish: compile, then copy this file to 2_Website/files/cv.tex and the
% PDF to 2_Website/files/cv.pdf.
%===============================================================================
\documentclass[letterpaper,11pt]{article}
\usepackage[utf8]{inputenc}
\usepackage{lmodern}
\usepackage[T1]{fontenc}
\usepackage[letterpaper, left=0.5in, right=0.5in, top=0.45in, bottom=0.55in]{geometry}

\usepackage{fancyhdr}
\pagestyle{fancy}
\fancyhf{}
\renewcommand{\headrulewidth}{0pt}
\renewcommand{\footrulewidth}{0pt}

\usepackage{enumitem}
\usepackage{titlesec}
\usepackage{xcolor}
\usepackage{textcomp}
\usepackage{hyperref}
\usepackage{array}

% Icons when fontawesome5 is installed (e.g. Overleaf); nothing otherwise, so a
% local build never prints stray bullets or "[in]".
\IfFileExists{fontawesome5.sty}{%
    \usepackage{fontawesome5}%
    \newcommand{\iconmail}{\faEnvelope\ }%
    \newcommand{\iconlinkedin}{\faLinkedin\ }%
    \newcommand{\iconweb}{\faGlobe\ }%
}{%
    \newcommand{\iconmail}{}%
    \newcommand{\iconlinkedin}{}%
    \newcommand{\iconweb}{}%
}

\definecolor{cvblue}{HTML}{1F5FA8}
\hypersetup{colorlinks=true, urlcolor=cvblue, linkcolor=cvblue,
            pdftitle={Guido Ardizzone -- CV}, pdfauthor={Guido Ardizzone}}

\setlength{\parindent}{0pt}
\setlength{\parskip}{0pt}

\titleformat{\section}
    {\large\bfseries\color{cvblue}}{}{0em}{}[{\color{cvblue}\titlerule[0.8pt]}]
\titlespacing*{\section}{0pt}{10pt}{6pt}

% Keep a section heading with its first lines: break the page early if fewer
% than ~5 lines remain (inline version of the needspace package).
\usepackage{etoolbox}
\makeatletter
\newcommand{\keepspace}[1]{\par\penalty-100\begingroup
    \setlength{\dimen@}{#1}\dimen@ii\pagegoal\advance\dimen@ii-\pagetotal
    \ifdim\dimen@>\dimen@ii\ifdim\dimen@ii>\z@\vfil\fi\break\fi\endgroup}
\makeatother
\preto\section{\keepspace{5\baselineskip}}

% Compact entry: institution / dates, then role / location in small italics
\newcommand{\cventry}[4]{%
    \textbf{#1}\hfill\textbf{#2}\\
    {\small\textit{#3}}\hfill{\small\textit{#4}}\par\vspace{5pt}}

% Paper: title line (bold title + coauthors/venue), then the abstract
\newcommand{\paper}[3]{%
    \textbf{#1}#2\par\vspace{4pt}%
    {\small\textbf{Abstract}: \textit{#3}}\par\vspace{10pt}}

% Press piece: linked title -- outlet, date
\newcommand{\press}[3]{%
    \href{#1}{#2} -- \textit{#3}\par\vspace{4pt}}

% Two-column rows (date | content)
\newenvironment{rows}{%
    \begin{tabular}{@{}p{2.8cm}p{\dimexpr\textwidth-2.8cm-2\tabcolsep}@{}}}%
    {\end{tabular}}

\begin{document}

%===============================================================================
\begin{center}
    {\Huge\scshape\textcolor{cvblue}{Guido Ardizzone}}\\[6pt]
    {\small Last Update: September 2026}\\[3pt]
    {\small
    \iconmail\href{mailto:guido.ardizzone@graduateinstitute.ch}{guido.ardizzone@graduateinstitute.ch}
    \quad\textperiodcentered\quad
    \iconlinkedin\href{https://www.linkedin.com/in/guido-ardizzone-b52048174/}{LinkedIn}
    \quad\textperiodcentered\quad
    \iconweb\href{https://guidoardizzone.github.io}{guidoardizzone.github.io}}
\end{center}
\vspace{-6pt}

%===============================================================================
\section*{Research Fields}

\textbf{Primary:} Sovereign Debt, Macro-Finance, Banking\\[2pt]
\textbf{Secondary:} International Finance, Sanctions and Geopolitical Risk

%===============================================================================
\section*{Education}

\cventry{Geneva Graduate Institute (IHEID)}{Sep. 2022 -- Present}{PhD in Economics}{Geneva, Switzerland}
\vspace{-4pt}{\small Supervisor: Ugo Panizza. Second reader: Rui Esteves.}\par\vspace{6pt}

\cventry{Bocconi University}{Sep. 2018 -- Apr. 2021}{MSc in Economics and Social Sciences, cum laude}{Milan, Italy}

\cventry{Bocconi University}{Sep. 2015 -- Oct. 2018}{BSc in Economics and Social Sciences}{Milan, Italy}

%===============================================================================
\section*{Work Experience}

\cventry{International Monetary Fund (IMF)}{Jun. 2026 -- Aug. 2026}{Fund Intern, Debt Policy Division (Strategy, Policy \& Review Department)}{Washington, DC, USA}

\cventry{The World Bank Group}{Jul. 2025 -- Aug. 2025}{Visiting PhD Student}{Washington, DC, USA}

\cventry{Geneva Graduate Institute (IHEID)}{Jan. 2023 -- Jun. 2025}{Research Assistant to Prof. Rui Esteves (SNF-funded project)}{Geneva, Switzerland}

\cventry{The World Bank Group}{Jan. 2023 -- Jun. 2023}{Short-Term Consultant (part-time, remote)}{Geneva, Switzerland}

\cventry{BAFFI CAREFIN Research Center, Bocconi University}{Sep. 2021 -- Jun. 2022}{Research Assistant (Carlo Altomonte)}{Milan, Italy}

\cventry{UCLA Anderson School of Management}{Jul. 2021 -- Jun. 2022}{Research Assistant (Paola Giuliano, remote)}{Milan, Italy}

\cventry{European Commission}{Feb. 2021 -- Jun. 2021}{Trainee, Chief Economist Team, DG Competition}{Brussels, Belgium}

%===============================================================================
\section*{Teaching Experience}

\cventry{Geneva Graduate Institute (IHEID)}{Sep. 2024 -- Present}{Teaching Assistant}{Geneva, Switzerland}
\vspace{-2pt}
\begin{rows}
Fall '24, '25   & Macroeconomics I (Master)\\[2pt]
Spring '25, '26 & Microeconomics II (Master)
\end{rows}

%===============================================================================
\section*{Grants and Scholarships}

\begin{rows}
2026       & Gallatin Excellence Scholarship (USD 20{,}000), visiting semester at the \href{https://psfl.princeton.edu/}{Princeton Sovereign Finance Lab}\\[3pt]
2022--2026 & IHEID Excellence Scholarship\\[3pt]
2021       & Zegna Founder's Scholarship Program, Ermenegildo Zegna Group (USD 5{,}000)
\end{rows}

%===============================================================================
% Research starts on its own page so abstracts are never split from their titles.
\clearpage
\section*{Working Papers and Work in Progress {\normalsize\mdseries(* Job Market Paper)}}
\vspace{2pt}

\paper{Sovereign Risk and Firm Debt Maturity*}{}{%
This paper studies how banks adjust the maturity of credit when sovereign stress rises.
I use confidential administrative data from the Banco de Portugal that track monthly
bank--firm credit relationships and the maturity of outstanding credit during the
2010--2012 European sovereign debt crisis. I exploit differences in banks' pre-crisis
exposure to distressed sovereign debt to compare lending to the same firm by banks with
different exposure. I find that banks do not reduce lending proportionately across
maturities. Banks more exposed to sovereign risk cut long-term credit while maintaining
short-term lending to the same firms, shifting the maturity of credit within continuing
relationships. The contraction in long-term credit is larger for small, highly leveraged,
and financially constrained firms. At the firm level, total bank credit also falls: firms
partially replace the long-term credit they lose with short-term borrowing, but
less-exposed lenders do not expand lending and firms are less likely to form new
relationships. Sovereign stress therefore affects both the quantity and maturity of bank
financing, leaving firms with less credit and shorter financing horizons.}

\paper{Senior Debt and Market Access}{ with Agust\'in Velasquez and Mahamoud Islam}{%
Multilateral lending is a significant component of sovereign debt in emerging and
developing economies. Because multilateral creditors receive de facto preferred creditor
treatment, a larger stock of multilateral claims can affect the pricing of private
sovereign debt through two opposing channels. On the one hand, official lending can ease
financing constraints and rollover pressure, reducing the likelihood of distress. On the
other hand, the seniority of multilateral claims leaves private creditors with a smaller
residual claim in states of distress, potentially increasing their required compensation.
We study how these forces interact using an annual panel of 55 emerging and developing
economies over 2004--2024. We measure senior creditor exposure as the share of external
public and publicly guaranteed debt owed to multilateral creditors and relate it to
sovereign spreads from the JP Morgan EMBI Global. Using a quadratic specification with
country and year fixed effects, we find that (i) on average, the relationship between
senior creditor exposure and sovereign spreads is non-linear, with the marginal
association becoming positive at about 40 percent of external debt; (ii) the
non-linearity is state dependent and concentrated among more vulnerable sovereigns; and
(iii) it is specific to multilateral creditor exposure, consistent with the role of de
facto seniority. The results are robust to alternative specifications, additional
controls, global risk factors, and excluding IMF credit from the exposure measure. We
interpret the evidence as consistent with a trade-off between the catalytic benefits of
official finance and the dilution of private claims created by senior creditor exposure.}

\paper{Sanctions, Lending, and the Bank Ownership Channel}{ with \href{https://sites.google.com/view/matteoficarra/home-page}{Matteo Ficarra}}{%
This paper examines the impact of sanctions on bank lending with a focus on the role of
bank ownership. Using balance sheet and ownership data for over 1,800 banks in 56
emerging market economies and 259 sanction episodes between 1995 and 2020, we find that
in response to sanctions foreign banks significantly reduce credit supply, particularly
to firms, in line with a flight-to-safety type of behaviour. As private domestic banks
are unable to increase lending, state-owned banks partially absorb credit demand by
expanding loans to both firms and households, suggesting a countercyclical response.
This asymmetric response is consistent across different sanction types, but particularly
strong for trade and financial sanctions. Our findings unveil a new channel to assess
the effectiveness of sanctions, highlighting the importance of banks' ownership
structure, and offers novel insights on the response of cross-border credit flows to
sanctions.}

%===============================================================================
% Pieces are bylined collectively to the Tortuga think tank (no individual authors).
\section*{Policy Publications and Press {\normalsize\mdseries(written with Tortuga)}}

\press{https://lavoce.info/archives/111182/campioni-europei-di-debito/}{Campioni europei di debito}{lavoce.info, 28/04/2026}
\press{https://lavoce.info/archives/102561/nadef-e-patto-di-stabilita-il-diavolo-e-nei-dettagli/}{Nadef e Patto di stabilità: il diavolo è nei dettagli}{lavoce.info, 24/10/2023}
\press{https://www.econopoly.ilsole24ore.com/2023/01/18/riforma-patto-stabilita-crescita/}{Perché la riforma del Patto di Stabilità e Crescita così com'è non va}{Econopoly (Il Sole 24 Ore), 18/01/2023}
\press{https://www.fanpage.it/politica/quanto-costa-realizzare-il-programma-elettorale-del-partito-democratico/}{Quanto costa realizzare il programma elettorale del Partito Democratico}{Fanpage.it, 14/09/2022}
\press{https://www.fanpage.it/politica/quanto-costa-realizzare-il-programma-elettorale-del-terzo-polo-di-azione-e-italia-viva/}{Quanto costa realizzare il programma elettorale del Terzo polo di Azione e Italia Viva}{Fanpage.it, 12/09/2022}

%===============================================================================
\section*{Presentations}

\begin{rows}
2026      & IMF Debt Policy Division Seminar (Washington, DC)\\[3pt]
2025--2026 & IHEID--UNIGE Joint PhD Workshop (Geneva), Claude Code for Economic Research Workshop (Geneva), IHEID Brown Bag Lunch (Geneva)
\end{rows}

%===============================================================================
\section*{Other Professional Services}

Seminar Organizer: \textit{IHEID Brown Bag Lunch PhD Seminars (2023--2026)}

%===============================================================================
\section*{Research Affiliations}

\cventry{\href{https://www.tortuga-econ.it/}{Tortuga}}{2021 -- 2026}{Senior Fellow}{}

%===============================================================================
\section*{Languages and IT Skills}

\begin{itemize}[leftmargin=1.2em, itemsep=1pt, topsep=0pt, label=\textbullet]
    \item \textbf{Languages}: Italian (native), English (fluent, IELTS 8.5), Spanish (good)
    \item \textbf{Software}: Stata, R, Dynare, \LaTeX, Claude Code, MS Office
\end{itemize}

%===============================================================================
\section*{References}

% Web-safe: referee email addresses removed (available on request).
\begin{tabular}{@{}p{0.48\textwidth}p{0.48\textwidth}@{}}
Prof. Ugo Panizza                                         & Prof. Rui Esteves\\
Professor of Economics and Pictet Chair in Finance and Development & Full Professor\\
Geneva Graduate Institute (IHEID)                         & Geneva Graduate Institute (IHEID)
\end{tabular}

\end{document}
